\section{Mechanistic Experiments and Diagnostics}
\label{app:mechanistic}

\subsection{Replay as a dynamical clue}
\label{app:replay-clue}

The temporal clue in Section~\ref{sec:sharpening} concerns the total online gap, not a constant increase in novel-continuation loss. In the seed-42 both-table run \code{s1v5_128_ep1xL4_20ep_both_fixed}, with the standard input-injection setting and 337 steps per pass, the gaps at the ends of passes 1--4 (steps 337, 674, 1011 and 1348) are $-0.062$, $1.199$, $2.730$ and $3.917$ nats. The first three replay increments are therefore $1.261$, $1.531$ and $1.187$ nats: comparable early increments, rather than an exact constant. The increments subsequently decline, reaching $0.241$ nats from pass 19 to 20. The essential observation is that the same missing continuations remain unobserved on every pass while the model's gap keeps growing. This motivates a dynamical account of the model response; it does not by itself distinguish optimization mechanisms.

\subsection{Sharpening diagnostics}

\begin{figure}[!t]
\centering
\begin{subfigure}[b]{0.58\linewidth}
\centering
\includegraphics[width=\linewidth]{figures/walk_their_dog_loss.pdf}
\subcaption{}
\label{fig:walk-their-dog}
\end{subfigure}\hfill
\begin{subfigure}[b]{0.38\linewidth}
\centering
\includegraphics[width=\linewidth]{figures/fig09_crowd_ce_seen_novel.pdf}
\subcaption{}
\label{fig:crowd-ce}
\end{subfigure}
\caption{Crowding-out: (a) context-only continuation losses for ``walk their dog''; (b) validation cross-entropy of train-seen and novel continuations, pooled over contexts.}
\label{fig:crowding}
\end{figure}

The diagnostic cited in Section~\ref{sec:sharpening} compares a context-only probe of ``walk their dog,'' whose training continuation is \emph{down} and whose fixed-validation continuations include \emph{daily} and \emph{for}, across replay passes (Figure~\ref{fig:walk-their-dog}). The aggregate diagnostic separates validation positions with a training-seen bigram context by whether their true continuation was also seen after that context in training, then compares the both-table model with the no-gram control (Figure~\ref{fig:crowd-ce}). Predictive entropy and cross-entropy are computed on the same fixed validation positions so that sharper distributions can be distinguished from improved accuracy.

For the displayed context, \emph{down} occurs once after the trigram in training; \emph{daily} and \emph{for} occur once each in the fixed validation window and not after this trigram in training. Panel (a) uses the seed-42 \path{sharpening_f99_contextonly_20260921_s42_fixed} context probe. At the ends of five passes, steps 337, 674, 1011, 1348, and 1685, it presents only the literal context and reads its full-vocabulary next-token distribution. Its training loss is $-\log p(\textit{down})$; the validation loss is the mean of $-\log p(\textit{daily})$ and $-\log p(\textit{for})$. Neither is the original full-prefix online loss. This one example illustrates crowding-out; it is not a claim that every context sharpens monotonically.

Panel (b) instead uses full causal prefixes and position-weighted validation cross-entropy. The archived aggregation \path{fig6_global_seen_novel_20260922_s42_fixed} plots the original both-table checkpoints and a separately evaluated no-gram control at the same five epoch ends. Of the 589{,}248 classified validation positions, 405{,}058 have a training-seen continuation and 158{,}790 a novel continuation after a training-seen bigram context; 25{,}400 contexts absent from training are excluded from these two curves. The entropy comparison also concerns fixed validation positions, not the context-only illustration. The two panels have different inputs and estimands and should not be read as a matched per-example loss decomposition.

\subsection{Context-loss masking}

The experiment in Section~\ref{sec:context-locality} excludes next-token loss for a random 2\% of training trigram contexts beginning at epoch 2 or 3, while retaining their input tokens, trainable parameters, and the loss denominator. A disjoint 2\% of unmasked contexts and a normal-training run from the same checkpoint are controls. Figure~\ref{fig:context-subspace-mask}(b) compares the masked and peer sets using $C_S=\sum_{i\in S}q_i[\bar\ell_{\mathrm{val},i}-\bar\ell_{\mathrm{tr},i}]$ on fixed batches. Figure~\ref{fig:context-subspace-mask}(a) tracks full-vocabulary softmax mass on training-seen and validation-only continuations, summed within each support and averaged over fixed validation positions. Validation-only continuations follow the context in fixed validation but not in training. These are group masses, not mean per-token probabilities; the two supports need not exhaust the vocabulary. The contrast tests whether removing a context's own updates stops its gap growth despite continued optimization on other contexts; it does not assume mathematically independent parameter subspaces.

The redrawn four-pass diagnostic preserves the archived values. Its mass panel uses seed-42 \code{s1v5_128_marm_both_ctxmask_e1_l4_fixed} and \code{s1v5_128_marm_nogram_ctxmask_e1_l4_fixed} at steps 337, 674, 1011 and 1348, with 8,167 masked and 8,266 unmasked-peer validation positions. The gap panel uses the both-table \code{ctxmask_ctrl_l4}, \code{ctxmask_e1_l4} and \code{ctxmask_e2_l4} runs with the same prefix and \code{_fixed} suffix, plotting their archived fixed-probe $C_A$ and $C_B$, not the online $C^{\mathrm{pos}}$ measurements. Run identifiers and plotted values are now traced; this does not resolve the protocol comparison below.

An independently archived one-pass audit makes the cohort and measurement choices explicit. It samples 2\% of exact training trigram contexts uniformly without replacement (selection seed 20260919); among 379{,}789 selected contexts, 5{,}565 have fixed-validation support (8{,}083 positions). A disjoint unmasked peer sample has 5{,}571 measurable contexts (8{,}082 positions), not the entire unmasked complement. In that audit, $q_i=N_{\mathrm{val},i}/589{,}824$ uses the original fixed-validation position count without renormalizing within either cohort. Per-context train and validation losses preserve their full causal prefixes at a common checkpoint, unlike the context-only sharpening probe. The epoch-2 pair is \path{ctxloss_e2_all_p02_control_s42_fixed} and \path{ctxloss_e2_all_p02_mask_s42_fixed} over steps 337--674; the separate epoch-3 pair is \path{ctxloss_all_p02_control_s42_fixed} and \path{ctxloss_all_p02_mask_s42_fixed} over steps 674--1011. Within each pair, normal and masked training begin from the same checkpoint, but these are separate experiments and are not the four-epoch trajectories displayed in Figure~\ref{fig:context-subspace-mask}(b). Suppressed aggregate growth on selected contexts alongside continued growth on peers supports approximate locality of additional gap accumulation. Individual losses still move in both directions, and a pre-existing gap can persist.

\subsection{Status of the context-masking comparison}
\label{app:mask-reconciliation}

The original trajectories in Figure~\ref{fig:context-subspace-mask} suggest that removing a context's own training loss arrests its gap growth. Registered attempts to reproduce this pattern under the minimal setting, however, did not recover the arrest with either context-type masking or table-row masking. In the seed-42 input both-table runs \code{s1v5_128_marm_both_rowmask_ctrl_fixed} and \code{s1v5_128_marm_both_rowmask_e1_fixed}, masking starts at step 338; at step 670 the masked cohort's online gap contribution is $0.0353$ nats, compared with $0.0247$ in the unmasked control. With masking from step 675 in \code{s1v5_128_marm_both_rowmask_e2_fixed}, the corresponding step-790 contributions are $0.0551$ versus $0.0492$ nats. These runs retain the standard input setting, $2^{20}$ rows per table, RMSProp table LR $128\times$ and backbone LR $0.0006$.

The rerun quantity is the cohort's validation-position share multiplied by the difference between its validation mean and its online training mean. It is distinct from the matched fixed-probe quantity used in the original figure description. The archived run configurations and treatment of the loss denominator still need reconciliation with the reruns. Thus these reruns are not a matched refutation of the displayed trajectories, but they prevent treating context independence as an established property of the current setting. Shared-backbone updates may continue to alter a masked context, and context-specific representations alone do not imply independent optimization.

\pending{Before submission, register the entropy probe and its aggregation script, and reconcile the archived multi-epoch masking protocol with the minimal-setting reruns. The separate one-pass audit above is not the provenance of Figure~\ref{fig:context-subspace-mask}.}


\subsection{Actual context-only output distributions}
\label{app:actual-logits}

Figures~\ref{fig:actual-petting}--\ref{fig:actual-there} compare the seed-42 runs \code{ep5ckpt_v5_nogram_s42_fixed} and \code{ep5ckpt_v5_both_s42_fixed} at steps 337 and 1685 (passes 1 and 5). Each input is BOS plus the three context tokens; probabilities use a full-vocabulary softmax over the checkpoint logits. The upper panels show probability and the lower panels show log probability. The named animals are single tokens under the training tokenizer. Others is the arithmetic mean of the 20 largest remaining token probabilities after excluding those named tokens, with the set reselected for each distribution; it is not remaining probability mass. In the lower panels it is the logarithm of this mean.

The selected \emph{petting your} and \emph{introducing your} contexts come from a probe cohort where dog follows the context in training and cat does not. They illustrate suppressed alternative continuations in the $n$-gram model, alongside improvement for the no-gram control. These checkpoint comparisons do not isolate the causal effect of training another context. The literal \emph{there is a} query is included to distinguish the conceptual example in Figure~\ref{fig:replay-transfer} from observations: both dog and cat follow it in training, and its measured trajectory need not match the schematic. These are context-only examples from one training seed, not natural-prefix population estimates.

\begin{figure}[!htbp]
\centering
\includegraphics[width=.94\linewidth]{figures/fig_logits_petting_your.pdf}
\caption{Output distributions after \emph{petting your}: pass 1 (light/open) and pass 5 (dark/filled), with shared scales within each row.}
\label{fig:actual-petting}
\end{figure}

\begin{figure}[!htbp]
\centering
\includegraphics[width=.94\linewidth]{figures/fig_logits_introducing_your.pdf}
\caption{Output distributions after \emph{introducing your}, using the conventions of Figure~\ref{fig:actual-petting}.}
\label{fig:actual-introducing}
\end{figure}

\begin{figure}[!htbp]
\centering
\includegraphics[width=.94\linewidth]{figures/fig_logits_there_is_a.pdf}
\caption{Output distributions after \emph{there is a}, where both dog and cat occur in training. Conventions follow Figure~\ref{fig:actual-petting}.}
\label{fig:actual-there}
\end{figure}
